%%
%% fall-lecture03.tex
%% 
%% Made by Alex Nelson <pqnelson@gmail.com>
%% Login   <alex@lisp>
%% 
%% Started on  2026-10-04T11:22:43-0700
%% Last update 2026-10-04T11:22:43-0700
%% 

\lecture[Continuum description]{}

\begin{node}
We will be studying integral form of equations of motion. This
involves a rather lengthy calculation (possibly several).
\end{node}

\begin{definition}
We define the \define{Control Volume} (often abbreviated as ``CV'') to
be a closed region in space with boundary kinematics specified (e.g.,
a region fixed over time, or moves in a specific way).
\end{definition}

\begin{definition}
A \define{System} is a connected material ``continuum''.
\end{definition}

\begin{node}[Setup and assumptions]
Consider a system described by a volume $V_{s}(t)$ with
boundary $\boundary V_{s}(t)$, which evolves over time. Consider also a control
volume $V_{c}$ with boundary $\boundary V_{c}$ assumed to be
stationary.

Now we assume that at time $t_{0}$ that
\begin{equation}
V_{c}=V_{s}(t_{0}),\quad\mbox{and}\quad\boundary V_{c}=\boundary V_{s}(t_{0}).
\end{equation}
That is to say, the system and control volume initially coincide as
the same region at initial time $t_{0}$.
\end{node}

\begin{node}
Let $B(t)$ be the amount of some quantity $b(\vec{x},t)$ per unit
mass for the system. Then we can describe the amount $B_{c}(t)$ contained in the
control volume (which will vary with respect to time if the system
flows into or out of the control volume) which satisfies
\begin{equation}
B_{c}(t_{0})=B(t_{0})
\end{equation}
and
\begin{equation}
B_{c}(t)=\int_{V_{c}}\rho b(\vec{x},t)\,\D^{3}x.
\end{equation}
Similarly, the quantity $B(t)$ for the system may be computed as:
\begin{equation}
B(t)=\int_{V_{s}(t)}\rho b(\vec{x},t)\,\D^{3}x.
\end{equation}
Then at time $t+\D t$, we can compute
\begin{equation}
B_{c}(t+\D t)=\int_{V_{c}}\rho b(\vec{x},t+\D t)\,\D^{3}x,
\end{equation}
and
\begin{equation}
B(t+\D t)=\int_{V_{s}(t+\D t)}\rho b(\vec{x},t + \D t)\,\D^{3}x.
\end{equation}
Then if we describe $B_{\text{in}}$ as the amount of $b$ contained in
$V_{s}(t+\D t)\setminus V_{c}$, and if $B_{\text{out}}$ describes the
amount of $b$ contained in $V_{c}\setminus V_{s}(t+\D t)$, then we
find
\begin{equation}
B_{c}(t+\D t)=B(t+\D t)+B_{\text{in}}-B_{\text{out}}
\end{equation}
by the conservation of mass. Then
\begin{equation}
B_{c}(t+\D t)-B_{c}(t)=B(t+\D t)-B(t)+B_{\text{in}}-B_{\text{out}},
\end{equation}
or
\begin{equation}\label{eq:fall:lec03:first-major-step}
B_{c}(t+\D t)-B_{c}(t)+B_{\text{out}}-B_{\text{in}}=B(t+\D t)-B(t).
\end{equation}
This is the first major step in our calculation.
\end{node}

\begin{node}
Now, suppose in 3-dimensions, we look at $\boundary V_{c}$ and
specifically at some small area $\D A$ on $\boundary V_{c}$ with
outward-pointing unit normal vector $\UnitNormalVec{n}$. Then the
velocity vector $\vec{u}$ (which has an angle $\alpha$ with $\UnitNormalVec{n}$) going out gives us
\begin{subequations}
  \begin{align}
\D B_{\text{out}} &= \begin{pmatrix}\mbox{amount of $b$
    that}\\\mbox{crosses $\D A$ in time $\D t$}
\end{pmatrix}\\
&= \rho b \underbrace{\vec{u}\cdot\UnitNormalVec{n}\,\D A\,\D
  t}_{\text{volume of cylinder}}|_{|\alpha|<\pi/2}
  \end{align}
\end{subequations}
The angle $\alpha$ between $\vec{u}$ and $\UnitNormalVec{n}$ must be
less than $|\alpha|<90^{\circ}$ for this to be outgoing. Similarly,
for incoming molecules, we must use the inward-pointing unit normal
vector (which, in our notation, is just $-\UnitNormalVec{n}$):
\begin{equation}
\D B_{\text{in}}=-\rho b \vec{u}\cdot\UnitNormalVec{n}\,\D A\,\D t|_{|\alpha|>\pi/2}
\end{equation}
where the leading minus sign comes from the fact that we are using the
inward-pointing unit normal vector. Then
\begin{subequations}
  \begin{align}
B_{\text{out}}-B_{\text{in}} &= \oint_{\boundary V_{c}}\D B_{\text{out}}-\D B_{\text{in}}\\
&=\oint_{\boundary V_{c}}\rho b \vec{u}\cdot\UnitNormalVec{n}\,\D A\,\D t|_{|\alpha|<\pi/2}-(-\rho b \vec{u}\cdot\UnitNormalVec{n})\,\D A\,\D t|_{|\alpha|>\pi/2}\\
&=\oint_{\boundary V_{c}}\rho b \vec{u}\cdot\UnitNormalVec{n}\,\D A\,\D t.
  \end{align}
\end{subequations}
This is the second major step in our calculation.
\end{node}

\begin{node}
We can now recall
Eq~\eqref{eq:fall:lec03:first-major-step}
\begin{equation}
\underbrace{B_{c}(t+\D t)-B_{c}(t)}_{=\D B_{c}(t)}+B_{\text{out}}-B_{\text{in}}=\underbrace{B(t+\D t)-B(t)}_{=\D B(t)}.
\end{equation}
We ``divide through by $\D t$'' (wink wink) and use the second major
result to get
\begin{equation}
\frac{D B}{D t}=\frac{\partial B_{c}}{\partial t}+\oint_{\boundary V_{c}}\rho b \vec{u}\cdot\UnitNormalVec{n}\,\D A
\end{equation}
where $D/D t$ is the ``material derivative'' which accounts for
advective contributions.
\end{node}

\begin{theorem}[Reynolds transport]
For a possibly moving control volume $V_{c}(t)$, we have
\begin{equation}
\frac{D}{D t}\int_{V_{s}(t)}\rho b\,\D V_{s} =
\frac{\D}{\D t}\int_{V_{c}(t)}\rho b\,\D V_{c} +
\oint_{\boundary V_{c}(t)}\rho b(\vec{u}-\vec{u}_{c})\cdot\UnitNormalVec{n}\,\D A_{c}
\end{equation}
where $\vec{u}_{c}$ is the velocity of $V_{c}$ the control volume,
or intuitively
\begin{equation}
  \begin{pmatrix}\mbox{Rate of change}\\
    \mbox{in $B$ for}\\
    \mbox{the system}
  \end{pmatrix}=
  \begin{pmatrix}\mbox{Rate of change}\\
    \mbox{in $B$ for}\\
    \mbox{the CV}
  \end{pmatrix} + 
  \begin{pmatrix}\mbox{Net flux of $B$}\\
    \mbox{across the control volume}
  \end{pmatrix}
\end{equation}
\end{theorem}

\begin{remark}
If $V_{c}$ is stationary (i.e., it is not moving in time),
then we have
\begin{equation}
\frac{\D}{\D t}\int_{V_{c}}\rho b\,\D V_{c}
=\int_{V_{c}}\frac{\partial(\rho b)}{\partial t}\,\D V_{c}.
\end{equation}
This should be obvious to everyone who remembers basic vector
calculus. 
\end{remark}

\begin{example}
Suppose we wanted to look at mass. Then we have $b=1$, so
\begin{equation}
B=\int_{V}\rho\,\D V
\end{equation}
is the mass of the system. Physics tells us that mass is conserved, so
\begin{equation}
\frac{D B}{D t}=0
\end{equation}
must be satisfied. This gives us the condition (which is \emph{always}
true at this level of approximation\footnote{For quantum field theory,
it's not true. But no one is concerned about fluids at that scale.})
\begin{equation}
0=\partial_{t}\int_{V_{c}}\rho\,\D V_{c}+\oint_{\boundary V_{c}}\rho\vec{u}\cdot\UnitNormalVec{n}\,\D A.
\end{equation}
\end{example}

\begin{node}
If we further assume a \define{Steady Flow} (so partial derivatives
with respect to time vanish), then we have mass conservation give us:
\begin{equation}
0=\oint_{\boundary V_{c}}\rho\vec{u}\cdot\UnitNormalVec{n}\,\D A.
\end{equation}
\end{node}

\begin{node}
A weaker assumption than stead flow is that we have an
\define{Incompressible Flow} (where the amount of fluid in a fixed
control volume is constant, so $\rho$ is effectively constant), we have
\begin{equation}
0=\oint_{\boundary V_{c}}\vec{u}\cdot\UnitNormalVec{n}\,\D A.
\end{equation}
\end{node}